\documentclass[border=10pt]{standalone}
\usepackage{tikz,ctex}
\usetikzlibrary{positioning}

\begin{document}

\begin{tikzpicture}[
    % 全局设置
    node distance = 1.6cm and 1.4cm,
    every node/.style = {
        draw,
        rounded corners,
        thick,
        align = center,
        font = \sffamily\small,
        minimum width = 2.0cm,
        minimum height = 0.9cm
    },
    % 样式定义
    centerStyle/.style = {
        fill = blue!45,
        text = white,
        font = \sffamily\bfseries\large,
        line width = 2pt,
        minimum size = 2.8cm,
        circle
    },
    branchTitle/.style = {
        font = \sffamily\bfseries,
        line width = 1.5pt
    },
    subNode/.style = {
        fill = white,
        font = \sffamily\small,
        minimum width = 2.2cm,
        minimum height = 0.8cm
    },
    leafNode/.style = {
        fill = white,
        font = \sffamily\footnotesize,
        minimum width = 1.9cm,
        minimum height = 0.7cm
    },
    arrowStyle/.style = {
        ->,
        thick,
        >=stealth,
        shorten >=2pt
    }
]

    % ================= 1. 中心节点 =================
    \node[centerStyle] (center) {9.3.学习曲线 \\ Learning Curves};

    % ================= 2. 上方分支：基本概念 (红) =================
    \node[branchTitle, fill=red!15, above=of center] (basic) {基本概念 \\ Basics};

    \node[subNode, above=of basic, xshift=-6cm] (train) {训练误差 \\ Training Error};
    \node[subNode, above=of basic, xshift=-2cm] (valid) {验证误差 \\ Validation Error};
    \node[subNode, above=of basic, xshift=2cm] (trade) {偏差--方差权衡 \\ Bias--Variance};
    \node[subNode, above=of basic, xshift=6cm] (iter) {迭代次数 \\ Iterations};

    % ================= 3. 右侧分支：早停法 (蓝) =================
    \node[branchTitle, fill=orange!15, right=of center] (early) {早停法 \\ Early Stopping};

    \node[subNode, right=of early, yshift=4cm] (stop) {停止准则 \\ Stop Criterion};
    \node[subNode, right=of early, yshift=2cm] (complex) {有效复杂度 \\ Effective Complexity};
    \node[subNode, right=of early, yshift=0cm] (weight) {类似权重衰减 \\ Like Weight Decay};
    \node[subNode, right=of early, yshift=-2cm] (path) {权重路径 \\ Weight Path};
    \node[subNode, right=of early, yshift=-4cm] (recip) {倒数关系 \\ $\tau/\eta \sim 1/\lambda$};

    % ================= 4. 下方分支：双下降 (绿) =================
    \node[branchTitle, fill=green!15, below=of center] (double) {双下降 \\ Double Descent};

    \node[subNode, below=of double, xshift=-6cm] (interp) {插值阈值 \\ Interpolation Threshold};
    \node[subNode, below=of double, xshift=-2cm] (modern) {现代深度学习 \\ Modern DL};
    \node[subNode, below=of double, xshift=2cm] (epoch) {训练轮数 \\ Epochs};
    \node[subNode, below=of double, xshift=6cm] (data) {数据量影响 \\ Data Size Effect};

    % ================= 5. 左侧分支：实践启示 (紫) =================
    \node[branchTitle, fill=purple!15, left=of center] (practice) {实践启示 \\ Practical Insights};

    \node[subNode, left=of practice, yshift=3cm] (bigger) {更大模型 \\ Bigger Models};
    \node[subNode, left=of practice, yshift=1cm] (implicit) {隐式偏置 \\ Implicit Bias};
    \node[subNode, left=of practice, yshift=-1cm] (sgt) {随机梯度下降 \\ SGD};
    \node[subNode, left=of practice, yshift=-3cm] (ironic) {更多数据 \\ More Data};

    % ================= 连线逻辑 =================
    % 中心 -> 四大分支标题
    \draw[arrowStyle, red] (center) -- (basic);
    \draw[arrowStyle, blue] (center) -- (early);
    \draw[arrowStyle, green!60!black] (center) -- (double);
    \draw[arrowStyle, purple] (center) -- (practice);

    % 分支标题 -> 子节点 (上方)
    \draw[arrowStyle, red!60] (basic) -- (train.south);
    \draw[arrowStyle, red!60] (basic) -- (valid.south);
    \draw[arrowStyle, red!60] (basic) -- (trade.south);
    \draw[arrowStyle, red!60] (basic) -- (iter.south);

    % 分支标题 -> 子节点 (右侧)
    \draw[arrowStyle, blue!60] (early) -- (stop.west);
    \draw[arrowStyle, blue!60] (early) -- (complex.west);
    \draw[arrowStyle, blue!60] (early) -- (weight.west);
    \draw[arrowStyle, blue!60] (early) -- (path.west);
    \draw[arrowStyle, blue!60] (early) -- (recip.west);

    % 分支标题 -> 子节点 (下方)
    \draw[arrowStyle, green!60!black] (double) -- (interp.north);
    \draw[arrowStyle, green!60!black] (double) -- (modern.north);
    \draw[arrowStyle, green!60!black] (double) -- (epoch.north);
    \draw[arrowStyle, green!60!black] (double) -- (data.north);

    % 分支标题 -> 子节点 (左侧)
    \draw[arrowStyle, purple!60] (practice) -- (bigger.east);
    \draw[arrowStyle, purple!60] (practice) -- (implicit.east);
    \draw[arrowStyle, purple!60] (practice) -- (sgt.east);
    \draw[arrowStyle, purple!60] (practice) -- (ironic.east);

\end{tikzpicture}

\end{document}